\note{3}{Sat 21 Oct 2023 21:31}{Martingale Theory}
\tags{Martingale, Stopping Time, Convergence Theorem}

\section{Martingale Theory}

\subsection{Basic Properties}

\begin{definition}[Martingale]
	\label{defn:martingale}
	A \textit{martingale} is a sequence of random variables $X_n$ such that
	\begin{enumerate}
		\item $E|X_n| < \infty $ 
		\item $X_n \in \mathcal{F}_n$ for some filtration $\mathcal{F}_n$
		\item $E(X_{n+1} | \mathcal{F}_n) = X_n$
	\end{enumerate}.
	If $E(X_{n+1}|\mathcal{F}_n) \le X_n$ instead, $X$ is said to be a \textit{supermartingale}, and if $E(X_{n+1}|\mathcal{F}_n) \ge X_n$, $X$ is said to be a  \textit{submartingale}.
\end{definition}
\begin{intuition}
	A martingale is a stochastic process that is somewhat stable, since each step has expected value zero. Supermartingales tend to decrease; submartingales tend to increase. We can view martingales to be a process that is ``monotone'' in the sense that it tends to go up or go down, while also subject to random perturbations.
\end{intuition}

Martingales are related to harmonic functions:
\begin{proposition}
	\label{prop:martingaleHarmonic}
	If $B(t)$ is a Markov process in $\mathbb{R}^d$ with radially symmetric distribution, then $f(B(t))$ is
	\begin{itemize}
		\item a martingale if $f$ is harmonic;
		\item a supermartingale if $f$ is superharmonic;
		\item a submartingale if $f$ is subharmonic.
	\end{itemize}
\end{proposition}

The behavior of a martingale can also be seen over a long range:
\begin{proposition}
	For $n > m$,
	\begin{itemize}
		\item If $X$ is a supermartingale, then $E(X_n | \mathcal{F}_m) \le X_m$.
		\item If $X$ is a submartingale, then $E(X_n | \mathcal{F}_m) \ge X_m$.
		\item If $X$ is a martingale, then $E(X_n | \mathcal{F}_m) = X_m$.
	\end{itemize}
\end{proposition}
\begin{proof}
	(1) and (2) are symmetric; (1) and (2) imply (3). To prove (1), induct on $n$, and use tower property of conditional expectation to complete the induction step.
\end{proof}

Passing through a convex function can bring subharmonicity:
\begin{theorem}
	\label{thm:convexSubmartingale}
	If either of the statements are true:
	\begin{enumerate}
		\item $X_n$ is a martingale and $\varphi$ is a convex function such that $E|\varphi(X_n)| < \infty $;
		\item $X_n$ is a submartingale and $\varphi$ is an increasing convex function such that $E|\varphi(X_n)| < \infty $;
	\end{enumerate}
	then $\varphi(X_n)$ is a submartingale.
\end{theorem}
\begin{proof}
	Use Jensen's inequality and the assumptions.
\end{proof}
\begin{remark}
	Also by using Jensen's inequality, we can show that if $\varphi$ is convex and $f$ harmonic, then $\varphi \circ f$ is subharmonic.
\end{remark}
\begin{example}
	As a consequences, 
	\begin{enumerate}
		\item If $X_n$ is a martingale (or nonnegative submartingale), then $|X_n|^p$ is a submartingale for $p \ge 1$.
		\item If $X_n$ is a submartingale, then $\exp(X_n)$ is a submartingale.
		\item If $X_n$ is a submartingale, then $(X_n - a)^+$ is a submartingale.
		\item If $X_n$ is a supermartingale, then $X_n \wedge a$ is a supermartingale.
	\end{enumerate}
\end{example}

More facts:
\begin{itemize}
	\item A ()martingale $X_n$ has natural filtration $\mathcal{G}_n = \sigma(X_k: k\le n)$, and $X_n$ is a ()martingale wrt its natural filtration. The natural filtration is the coarsest filtration that makes $X_n$ a ()martingale.
	\item If $X_n$ and $Y_n$ are submartingales, then $X_n \vee Y_n$ is a submartingale.
	\item (Switching Principle) If $X_n$ and $Y_n$ are supermartingales, $N$ is stopping time so that $X_N > Y_N$ a.s. Then $Z_n := X_n 1_{N > n} + Y_n 1_{N \le  n}$ is a supermartingale.
\end{itemize}


There are three important examples of martingales:
\begin{itemize}
	\item (Linear martingale)
	Let  $X_1, X_2, \ldots$ be iid, and $S_n = \sum_{j = 1}^n X_j$. Then $S_n$ is a martingale adapted to $\mathcal{F}_n = \sigma(X_1, \ldots, X_n)$,  $\mathcal{F}_0$ the trivial $\sigma$-algebra. If we define $S_0 = a$, this is called a martingale starts at $a$.
\item (Quadratic Martingale)
	Let $X_ k$ be iid, $\mu = E(X_k) = 0$, $\sigma^2 = \operatorname{Var} X_k = E X_k^2$. Then $S_n^2 - n \sigma^2$ is a martingale.
\item (Exponential Martingale)
	If $X_k \ge 0$ and $E X_k = 1$ are iid, then $\prod_{k = 1}^n X_k $ is a martingale wrt the natural filtration.

	Now, if we take $X_k = \exp(\lambda \xi_k)$, and $\varphi(\lambda) = E \exp(\lambda \xi_k)$, then $M_n = \frac{1}{\varphi(\lambda)^n} \exp(\lambda S_n)$ is a martingale.
\end{itemize}

\begin{example}
	(Dyadic martingale)

	Let $\mathcal{F}_n$ be the dyadic filtration on $[0,1)$, defined as $\mathcal{F}_0 = \{\emptyset , [0,1)\} $, $\mathcal{F}_1 = \sigma([0,\frac{1}{2}), [\frac{1}{2}, 1))$, etc. Then for any random variable $X$, $E(X | \mathcal{F}_n)$ is a martingale.

	This can be directly checked from definition, but let's give it more detail. Each $X_n$ is a r.v. taking discrete values and only depend on which dyadic interval $I^n_k$ the ``true answer'' $x$ lies in. The distribution is given by $p_{X_n}(I^n_k) = \int _{I^n_k} f$ for $k = 1, \ldots, n$, provided that $X$ has a pdf $f$ on $[0,1]$.
\end{example}

\subsubsection{No-strategy theorem}

\begin{theorem}
	Applying a predictable sequence $H_n \ge 0$ to a submartingale will produce a submartingale. That is, if $H_n$ is predictable and $X_n$ is a submartingale, then $(H\cdot X)_n = \sum_{k \le n}H_k (X_k - X_{k-1})$ is a submartingale.
\end{theorem}
Same applies to supermartingales and martingales. For martingales, we don't need to assume $H_n \ge 0$.
\begin{proof}
	We actually need additional assumption that $E\left| H_n (X_n - H_{n-1}) \right| < \infty $ for each $n$, in order to justify using the identity $E\left( XY | X \right) = X E(Y|X)$. This assumption is usually easily satisfied, for example, by having bounded $H_n$ for each $n$.

	Direct computation:
	\[
		E\left( (H \cdot X)_{n+1} | \mathcal{F}_n \right) 
		= (H \cdot X)_n + H_{n+1} E\left( X_{n+1} - X_n | \mathcal{F}_n \right) 
		\le (H \cdot X)_n
	.\] 
\end{proof}






\subsection{Martingale Convergence Theorem}

We have analogue for monotone convergence theorem for sequences:

\begin{theorem}[Martingale Convergence Theorem]
	(Doob's Martingale Convergence Theorem)
	\label{thm:DoobMartingaleConvergence}
	If $X_n$ is a submartingale with $\sup EX_n^+ < \infty $, then $X_n$ converges a.s. to a limit $X$ with $E|X| < \infty $.
\end{theorem}
The proof uses ``upcrossing inequality''. The intuition is that, a martingale that doesn't converge needs to oscillate infinitely across some positive distance; $\lim\inf X_n < \lim\sup X_n$. But this creates too much a burden for the martingale. Martingales are ``lazy'' about repeatedly crossing the same interval up and down, characterized by the upcrossing inequality.
We give the proof in the section for inequalities.

\begin{corollary}
	If $X_n \ge 0$ is a supermartingale, then $X_n \to  X$ a.s. and $EX \le  E X_0$.
\end{corollary}




\subsection{Stopped Martingale}

\begin{definition}[Stopping Time]
	A stopping time $T: \Omega \to \mathbb{N}$ is defined by the property
	$\{T \le n\} \in \mathcal{F}_n, \quad \forall n$.
	Equivalently we can use $\{T = n\} $ in the definition.
\end{definition}

\begin{proposition}
	\begin{enumerate}
		\item 
	Fixed time $n$ is a stopping time.
\item If $T_1$ and $T_2$ are stopping times, so is $T_1 \wedge T_2$ and $T_1 \vee T_2$.
\item If $T_n$ is nondecreasing, then $T = \sup _n T_n$ is a stopping time.
\item If $T_n$ is nonincreasing, then $T = \inf _n T_n$ is a stopping time.
\item If $T$ is a stopping time, then $T + n$ is a stopping time for $n > 0$.
	\end{enumerate}
\end{proposition}

Stopping a martingale does not destroy its martingale property. Besides, the expected value of a martingale also remains stable:
\begin{theorem}
	\label{thm:EXTn}
	If $X_n$ is a submartingale, then $X_{T \wedge n}$ is a submartingale, and 
	\[
		E X_0 \le  E X_{T \wedge n}  \le  E X_n
	\]
	for all $n$. 

	Analogous results are true for supermartingales and martingales.
\end{theorem}
\begin{proof}
	We can define a predictable sequence $H_n = 1_{T \ge n}$ so that 
	$X_{T \wedge n} = \left( H \cdot X \right) _n$,
	and the no-strategy theorem tells us $(H \cdot X)_n$ is a submartingale, thus
	\[
		E X_0 = E (H \cdot X)_0 \le E(H \cdot X)_n = E X_{T \wedge n}
	.\] 
	For the second inequality, apply the same argument to $X_n - X_{T \wedge n} = (K * X)_n$ where $K_n = 1_{N \ge n}$.
\end{proof}


\subsubsection{Optional Stopping Theorem}

Optional Stopping Theorems refer to some results that control $E X_T$, in the same manner of Theorem~\ref{thm:EXTn}. Although it seems simple, it's actually not easy at all to remove this $n$.

Firstly, if $T$ gets very large, we might not have a very good control of $X_T$ so that $E X_T$ may be unbounded or even not exist; moreover, we have to handle the event $\{T = \infty \}$. To handle these issues, we require a condition of \textit{uniform integrability}.

We want to start from Theorem~\ref{thm:EXTn} and take limit $n \to \infty $. In doing so, we expect $X_{T \wedge n} \to  X_T$ both a.s. and in $L^1$, so that $X_T$ exists and $E X_T$ gets controlled. Those are exactly the conditions given by uniform integrability.

The following version of OST does not use uniform integrability, but assumes some extra conditions.
\begin{theorem}[Optional Stopping Theorem]
	Let $T$ be a stopping time and $X_n$ a martingale. Then $E X_T = E X_0$ if one of following is true:
	\begin{enumerate}
		\item $T$ is a.s. bounded by some deterministic $K$.
		\item $X_n$ is a.s. bounded by some deterministic $K$, and $T$ is finite a.s.
		\item $E(T) < \infty $ and $X_k$ has deterministically bounded increment.
	\end{enumerate}

	Same is true for submartingales and supermartingales, though we would have $E X_T \ge E X_0$ and $E X_T \le E X_0$.
\end{theorem}
\begin{proof}
	Condition (1) is very strong, so the proof is simple. Since we already know $E(X_{T \wedge n}) = E(X_0)$, we can take $n = K$ and obtain $EX_T = EX_0$.

	Under condition (2), $T$ finite implies $X_{T \wedge n} \to  X_T$ a.s., and boundedness of $X_n$ allows us to apply dominated convergence theorem, so $E X_{T \wedge n} \to  E X_T$, but $E X_{T \wedge n} \equiv E X_0$.

	Under condition (3), we have $|X_{T \wedge n} - X_0| \le  K T$, so $E|X_{T \wedge n}| < K ET < \infty $. Apply DCT and take $n \to \infty $, we have $E X_T = E X_0$.
\end{proof}

\begin{remark}
	For non-negative supermartingales, in view of Theorem~\ref{thm:DoobMartingaleConvergence}, we always have convergence  $X_n \to X_\infty $; and $0 \le X_n \le X_0$ gives us $E X_{T \wedge n} \le  E X_0$ by Theorem~\ref{thm:EXTn}. Apply Fatou's Lemma to get $E X_T \le  E X_0$.
\end{remark}

Under a much nicer condition of \textit{uniform integrability}, we don't need any extra assumptions:
\begin{theorem}
	if $X_n$ is a uniformly integrable submartingale, then for any stopping time $T$ (possibly infinite), we have 
	\[
		E X_0 \le  E X_T \le E X_\infty
	\] 
	where $X_\infty  = \lim _{n\to \infty } X_n$.
\end{theorem}

\begin{proof}
	First of all, for submartingales we have
	\[
		E X_0 \le  E(X_{n \wedge T}) \le  E X_n
	.\] 
	Since $X_n$ is a uniformly integrable submartingale, $X_{n \wedge T}$ is also (Proposition~\ref{prop:UIMartingale}). Now uniform integrability means $X_n$ converges to $X_\infty $ both a.s. and in $L^1$, and $X_{n \wedge T}$ converges to $X_T$ both a.s. and in $L^1$. Take limit in the inequality above.
\end{proof}

\subsubsection{Results about uniform integrability}

\begin{definition}[Uniform Integrability]
	A collection of random variables $X_i, i \in I$ is said to be \textit{uniformly integrable} if
	\[
		\lim_{M \to \infty } \left( \sup _{i \in I} E\left( |X_i| ; |X_i| > M \right)  \right)  = 0
	.\] 
\end{definition}

\begin{theorem}
	(Durett 4.6.3)
	Suppose $X_n$ is $L^1$ bounded for all $n$. If $X_n \to  X$ in probability, then the following are equivalent: 
	\begin{enumerate}
		\item  $\{X_n\} $ is uniformly integrable.
		\item $X_n \to  X$ in $L^1$
		\item  $E|X_n| \to E|X| < \infty $.
	\end{enumerate}
\end{theorem}
\begin{proof}
	We give a shortened version. First, (ii) and (iii) are equivalent: (ii) implies (iii) by Jensen's and triangle inequality; (iii) implies (ii) by the i.p. convergence: $E|X_n - X| \le E\left( |X_n - X|; > \varepsilon \right) + \varepsilon $. The first probability can be controlled since (iii) provides an upper bound for $|X_n - X|$.

	Now (i) implies (ii) can be seen from
	 \[
		E|X_n - X| \le E\left( |X_n - X|; > M \right) 
		+ E\left( |X_n - X|; < M; > \varepsilon \right) 
		+ E\left( |X_n - X|; < \varepsilon \right) 
	.\] 
	The first term vanishes by uniform integrability of $X_n - X$. The second term vanishes by i.p. convergence and upper bound $M$. The second term bounded by $\varepsilon$.

	(ii) implies (i): The trick is again to use vertical truncation on $X_n$ and $X$. From i.p. convergence  $|X_n| 1\left( |X_n| \le M \right) \to |X| 1\left( |X| \le M \right) $ and DCT, we have $E\left( |X_n|; |X_n| \le M \right) \to E\left( |X|; |X| \le M \right) $. Now for all $M$,
	\begin{align*}
		E\left( |X_n|; |X_n| \ge M \right) 
		&= E|X_n| - E\left( |X_n|; |X_n| \le M \right)  \\
		&\to E|X| - E\left( |X|; |X| \le M \right) &(n \to \infty ) \\
		&\to 0 &(M \to \infty )
	.\end{align*}

	Now for all $n > n_0(\varepsilon)$, we have $E(|X_n|; \ge M) < E\left( |X|; \ge M \right) + \varepsilon$. For large enough $M$, $E(|X_n|; \ge M) < 2 \varepsilon$. Increase $M$ to have $E(|X_n|; \le M) < 2 \varepsilon$ for $n \le n_0$. Hence $E\left( |X_n|; \ge M \right) \to 0$ uniformly as $M \to \infty $.

\end{proof}


\begin{theorem}
	For a submartingale the following are equivalent:
	\begin{enumerate}
		\item It is uniformly integrable
		\item It converges a.s. and in $L^1$
		\item It converges in $L^1$ 
		\item (For $X$ a martingale) there exists $X \in L^1$ such that $X_n = E(X | \mathcal{F}_n)$.
	\end{enumerate}
\end{theorem}
\begin{proof}
	First assume $X_n$ is a submartingale and prove $\text{(1)} \implies \text{(2)} \implies \text{(3)} \implies \text{(1)}$.

	(1) implies (2). Uniform integrability gives us $L^1 $ boundedness. Martingale convergence theorem gives us a.s. convergence, which implies convergence in probability (i.p.). Now the previous theorem gives us $L^1$ convergence.

	(2) implies (3). Trivial.

	(3) implies (1). $L^1$ convergence implies i.p. convergence.  It also implies $L^1$ boundedness. Now apply $\text{(2)}\implies\text{(1)}$ in the previous theorem.

	Now assume $X_n$ is a martingale, we prove $\text{(3)}\implies \text{(4)} \implies \text{(1)}$.

	(3) implies (4). For $n > m$, we have $X_m = E\left( X_n | \mathcal{F}_m \right)  \to  E\left( X | \mathcal{F}_m \right) $.

	(4) implies (1). We verify that any collection of conditional expectations of a $L^1$ random variable $X$ is uniformly integrable (not only for filtrations). That is, we want to show that given $X \in L^1$, $\{E(X | \mathcal{F}): \mathcal{F} \subset \mathcal{F}_0\} $ is uniformly integrable.

	It suffices to show that for any $\varepsilon$, there exists $M$ such that $E\left(|E(X | \mathcal{F})| ; > M\right) < \varepsilon$. By Jensen's inequality,
	\[
		E\left( |E(X|\mathcal{F})|;|E(X|\mathcal{F})|>M \right) 
		\le E\left( E(|X| |\mathcal{F}); |E(X|\mathcal{F})|>M \right) 
		= E(|X| ; |E(X|\mathcal{F})|> M)
	.\] 
	But $L^1$ convergence ensures that this goes to $0$ as $M \to  \infty $. First, DCT tells us that if $P(A) \to  0$, $E(|X|; A) \to 0$. Next, $P(|E(X | \mathcal{F})| > M) \le \frac{E|X|}{M} \to 0$ as $M \to  \infty $. 
	(The key moment is control of $|E(X|\mathcal{F})|$ by $E|X|$ and $P(|E(X|\mathcal{F})>M)$ by $E|X|/M$, both independent of $\mathcal{F}$, thus providing uniform bound.)
\end{proof}



\begin{proposition}
	(Conditions for uniform integrability)
	\label{prop:UICond}
	\begin{itemize}
		\item If $|X_i| \le  Y$ where $E Y < \infty $, then $X_i$ is a uniformly integrable family.
		\item If $\varphi\ge 0$,  $\varphi(x) / x \to  \infty $ as $x \to  \infty $, and $E \varphi(|X_i|) \le  C$ for all $i \in I$, then $X_i$ is uniformly integrable.
		\item If $X_n$ is a submartingale with bounded increments, and $N$ is a stopping time with $EN < \infty $, then $X_{N \wedge n}$ is uniformly integrable and $E X_N \ge  E X_0$.
	\end{itemize}
\end{proposition}

\begin{proposition}
	(Uniformly integrable martingales)
	\label{prop:UIMartingale}
	\begin{itemize}
		\item If $X_n$ is a uniformly integrable submartingale, then for any stopping time $T$, $X_{T \wedge n}$ is a uniformly integrable submartingale.
		\item The same applies to supermartingales and martingales.
		\item For any stopping time $T$, $X_n \mapsto X_{T \wedge n}$ preserves uniform integrability, for any martingale, submartingale or supermartingale.
	\end{itemize}
\end{proposition}
\begin{intuition}
	A stopped martingale $X_{T \wedge n}$, when plotted wrt time $n$, looks like a random process that remains at a fixed value after some time. $X(\omega)$ and $T(\omega)$ are not independent, for example, if  $T$ is hitting time of $a$, then $X_T = a$ on the set $T < \infty $.

	It appears at first that even if  $X_n$ is uniformly integrable, $X_T$ can be something much bigger, since we can let the process stop when we have reached high enough $M$. However, note that as we take $M \to \infty $ in the definition of U.I., $T$ does not depend on $M$.
\end{intuition}
\begin{proof}
	\begin{align*}
		E\left( |X_{T \wedge n}|; > K \right) 
		&= E\left( |X_{T \wedge n}|; > K, T \le n \right) 
		+ E\left( |X_{T \wedge n}|; > K, T > n \right) \\
		&= E\left( |X_{T}|; |X_{T}|> K, T \le n \right) 
		+ E\left( |X_{n}|; |X_n|> K, T > n \right) \\
		&\le E\left( |X_T|; > K \right) + E\left( |X_n|; > K \right) 
	.\end{align*} 
	Uniform integrability of $X_n$ tells us the second term goes to zero uniformly in $n$ as $K \to  \infty $. Now we want to show that $E|X_T|< \infty $. 
	
	Since $X_n$ is a submartingale, we have $E X_{T \wedge n}^+ \le  E X_n^+$. Hence
	\[
		\sup_n E X_{T \wedge n}^+ 
		\le 
		\sup_n E X_{n}^+ 
		< \infty 
	.\] 
	The martingale convergence theorem now tells us that the pointwise limit $X_{T \wedge n} \to  X_T$ is $L^1$ integrable, so $E|X_T| < \infty $.

	Here we could have just used the DCT to see that $X_T$ is $L^1$. (check)
\end{proof}






\subsection{Martingale Inequalities}
Next: Lp convergence, (concentration) inequalities, upcrossing inequality, tail sigma fields and  01 laws, square integrability









